% part: intuitionistic-logic
% chapter: propositions-as-types
% section: proofs-to-terms

\documentclass[../../../include/open-logic-section]{subfiles}

\begin{document}

\olfileid{int}{pty}{pt}

\olsection{\usetoken{P}{derivation}ને સાબિતી-પદોમાં રૂપાંતરિત કરવી}

સ્વાભાવિક નિગમનની !!{proof}sને સાબિતી-પદોમાં
રૂપાંતરિત કરવી, એટલે~\Log{N2i}ની !!{proof}sને
\Log{tN2i}ની !!{proof}sમાં ફેરવવી, ખૂબ
સરળ છે. !!{derivation}ના દરેક સિક્વન્ટની
જમણી બાજુના !!{formula}ની પહેલાં
સાબિતી-પદ લખવાનું જ છે; આમ
$\Gamma \Sequent M : !A$ સ્વરૂપના
સિક્વન્ટો મળે છે. પછી કહીએ કે
સંદર્ભ~$\Gamma$માં $M$ એ~$!A$નું
\emph{સાક્ષ્ય આપે છે}. કયું સાબિતી-પદ
લખવું તે \Log{tN2i}ના નિયમોથી
સંપૂર્ણ નક્કી થાય છે.

\begin{prop}
\Log{N2i}માં $\Gamma \Sequent !A$ની
દરેક !!{proof}~$\delta$ને અનુરૂપ
\Log{tN2i}માં $\Gamma \Sequent N : !A$ની
!!a{proof}~$\delta'$ છે. સાબિતી-પદ~$N$
એ~$\delta$થી અનન્ય રીતે નક્કી થાય છે.
\end{prop}

\begin{proof}
$\delta$ની ઊંચાઈ પર આગમન કરીએ.
$\delta$ એ $x:!A \Sequent !A$ સ્વયંસિદ્ધ હોય,
તો $\delta'$ એ~$x:!A \Sequent x:!A$
સ્વયંસિદ્ધ છે. $\delta$ અનુમાન-નિયમથી
પૂરી થતી હોય, તો આગમન-ધારણા દરેક
આધારવિધાનનું સાબિતી-પદ અને તેને ધરાવતા
અનુરૂપ સિક્વન્ટની \Log{tN2i}-!!{proof}s
આપે છે. અનુરૂપ \Log{tN2i} નિયમ
લાગુ કરીએ; આ \Log{tN2i} અનુમાન-નિયમ સંબંધિત
સાબિતી-પદ નક્કી કરે છે.
\end{proof}

\begin{ex}
  \Log{N1i}માં $(!A \land !B) \lif (!B \land !A)$ની
  આ ખૂબ સરળ !!{proof} લો:
  \[
  \AxiomC{$\Discharge{!A\land !B}{x}$}
  \RightLabel{\Elim\land[2]}
  \UnaryInfC{$!B$}
  \AxiomC{$\Discharge{!A\land !B}{x}$}
  \RightLabel{\Elim\land[1]}
  \UnaryInfC{$!A$}
  \RightLabel{\Intro\land}
  \BinaryInfC{$!B \land !A$}
  \DischargeRule{\Intro\lif}{x}
  \UnaryInfC{$(!A \land !B) \lif (!B \land !A)$}
  \DisplayProof
  \]
  \Log{N2i}માં તે આવી છે:
  \[
  \Axiom$x : !A\land !B \fCenter !A\land !B$
  \RightLabel{\Elim\land[2]}
  \UnaryInf$x : !A\land !B \fCenter !B$
  \Axiom$x : !A\land !B \fCenter !A\land !B$
  \RightLabel{\Elim\land[1]}
  \UnaryInf$x : !A\land !B \fCenter !A$
  \RightLabel{\Intro\land}
  \BinaryInf$x : !A\land !B \fCenter !B \land !A$
  \RightLabel{\Intro\lif}
  \UnaryInf$\fCenter (!A \land !B) \lif (!B \land !A)$
  \DisplayProof
  \]
  ઉપરથી નીચે સાબિતી-પદો સોંપીએ.
  સ્વયંસિદ્ધોનાં પદો સંદર્ભના ચલ~$x$
  જ છે. પછી $\Elim\land[i]$ને અનુરૂપ
  પદો $\proj{2}{x}$ અને~$\proj{1}{x}$
  તરીકે નક્કી થાય છે. \Intro{\land}
  લાગુ કરીએ ત્યારે આ બંને પદો જોડીએ:
  $\pair{\proj{2}{x}}{\proj{1}{x}}$.
  છેલ્લે \Intro{\lif} નિયમ
  $\lambd$-અમૂર્તીકરણ લાગુ કરે છે:
  તે~$x$ને બદ્ધ કરે છે અને~$x$ એ
  ધારણા~$!A \land !B$નું ચિહ્ન હતું
  તે નોંધે છે:
  $\lambd[\typeof{x}{!A \land
  !B}][\pair{\proj{2}{x}}{\proj{1}{x}}]$.
  આમ \Log{tN2i}ની !!{proof} છે:
  \[
  \Axiom$x : !A\land !B \fCenter x: !A\land !B$
  \RightLabel{\Elim\land[2]}
  \UnaryInf$x : !A\land !B \fCenter \proj{2}{x} : !B$
  \Axiom$x : !A\land !B \fCenter x: !A\land !B$
  \RightLabel{\Elim\land[1]}
  \UnaryInf$x : !A\land !B \fCenter \proj{1}{x} : !A$
  \RightLabel{\Intro\land}
  \BinaryInf$x : !A\land !B \fCenter \pair{\proj{2}{x}}{\proj{1}{x}} : 
    !B \land !A$
  \RightLabel{\Intro\lif}
  \UnaryInf$\fCenter \lambd[\typeof{x}{!A \land
  !B}][\pair{\proj{2}{x}}{\proj{1}{x}}] : (!A \land !B) \lif (!B \land !A)$
  \DisplayProof
  \]
\end{ex}


\end{document}
